\chapter{Континуальная внутренняя щель полного прямого вихря}
\label{ch:v6-bosonic-defect-full-tensor-internal-gap-gate}

\section{Постановка}

Высокоугловой гейт закрыл знак полного поперечного оператора для всех
скрученных характеров и угловых меток. Оставалось вычислить точный масштаб
первого внутреннего возбуждения после удаления переносной пары. В
спектральной теории вихрей именно отделённая от нуля положительная мода
задаёт первый локализованный внутренний канал
\cite{full-tensor-gap-alonso}.

На сетке $80$ глобальный кандидат находился в сопряжённых блоках
\begin{equation}
 (c,j)=(1,1),\qquad (-1,-1).
\end{equation}
Для проверки его глобальности одновременно отслеживались ближайшие блоки
$(\pm1,0)$, прежняя ошибочная пара $(1,-1),(-1,1)$ и нейтральный
осесимметричный сектор.

\section{Сгущение критического блока}

На сетках
\begin{equation}
 N=70,100,140,200,280,400,560
\end{equation}
нижний уровень критической пары равен
\begin{equation}
 3.61859362,
 3.61897799,
 3.61915522,
 3.61924813,
 3.61929150,
 3.61931439,
 3.61932514.
\end{equation}
Максимальное различие сопряжённых блоков не превышает
$3.29\cdot10^{-14}$, а максимальный эрмитов остаток равен
$8.39\cdot10^{-14}$.

Линейная экстраполяция последних четырёх сеток по $1/(N-1)^2$ даёт
\begin{equation}
 \Delta_{\mathrm{int}}=3.6193363303.
\end{equation}
Независимый трёхпараметрический фит
$\Delta_N=\Delta_\infty+A/(N-1)^p$ даёт
\begin{equation}
 \Delta_\infty=3.6193362312,
 \qquad p=2.00564.
\end{equation}
Различие двух пределов меньше $10^{-7}$; сдвиг между сетками $400$ и $560$
равен $1.07\cdot10^{-5}$.

\section{Изоляция первого внутреннего уровня}

Ближайшая конкурирующая пара $(c,j)=(1,0),(-1,0)$ имеет предел
\begin{equation}
 \Delta_{\mathrm{next}}=3.6649900191,
\end{equation}
поэтому отделение равно
\begin{equation}
 \Delta_{\mathrm{next}}-\Delta_{\mathrm{int}}
 =0.0456536888>0.
\end{equation}
Второй уровень того же критического блока стремится к $3.6869379056$, а
прежний отрицательный кандидат после калибровки --- к $3.6898947863$.
Следовательно, найденная пара является изолированным глобальным минимумом
внутреннего спектра.

\begin{theorem}
После факторизации двух переносных нулей полный поперечный гессиан
стационарного прямого вихря имеет строго положительную континуальную щель
\begin{equation}
 \Delta_{\mathrm{int}}=3.61933633.
\end{equation}
Прямая нить линейно устойчива относительно всех поперечных тензорных и
калибровочных возмущений.
\end{theorem}

\begin{nogocriterion}
Положительная щель прямой нити не доказывает существование частицы. При
изгибе появляются кривизна мировой линии, взаимодействие удалённых участков
и глобальные моды радиуса петли. Эти члены отсутствуют в прямом поперечном
операторе и должны быть выведены из того же родительского действия.
\end{nogocriterion}

\section{Следующий различающий тест}

Нулевая гипотеза: медленно изогнутая нить наследует действие Намбу--Гото с
натяжением $1.5744526907$ и конечную поправку жёсткости, вычислимую через
разрешённый поперечный гессиан со щелью $3.61933633$.

Стоп-критерий: если коэффициент кривизны не определяется существующим
родительским действием или делает энергию кольца монотонной по радиусу,
прямая устойчивая нить не замыкается в локализованную частицу без нового
физического члена.

\begin{protocol}
\begin{enumerate}
 \item переносная пара: \textbf{удалена};
 \item критические характеры: \textbf{$c=\pm1$};
 \item критические метки: \textbf{$j=\pm1$};
 \item континуальная щель: \textbf{$3.61933633$};
 \item порядок сходимости: \textbf{$2.00564$};
 \item отрыв от конкурента: \textbf{$0.04565369$};
 \item полная поперечная устойчивость прямой нити: \textbf{закрыта};
 \item эффективное действие кривой нити: \textbf{открыто};
 \item замкнутая петля: \textbf{не проверена};
 \item рождение материи: \textbf{не закрыто}.
\end{enumerate}
\end{protocol}

Машинный сертификат сохранён в
\path{s2t_v6_bosonic_defect_full_tensor_internal_gap_gate_results.json}.

\begin{thebibliography}{9}
\bibitem{full-tensor-gap-alonso}
J.~M.~Alonso-Izquierdo, W.~Garc\'ia Fuertes, J.~Mateos Guilarte,
\emph{Dissecting zero modes and bound states on BPS vortices in
Ginzburg--Landau superconductors}, arXiv:1602.09084.
\end{thebibliography}